\section*{Tóm tắt}
\addcontentsline{toc}{section}{Tóm tắt}

Đồ án khảo sát hai mô hình Transformer tiền huấn luyện cho bài toán \textbf{đọc hiểu trích xuất tiếng
Việt} trên UIT-ViQuAD~2.0: \textbf{PhoBERT} --- đơn ngữ, nhận đầu vào \emph{đã tách từ} --- và
\textbf{XLM-R} --- đa ngữ, dùng subword trên văn bản thô. Hai mô hình khác nhau đúng ở khâu tách từ, nên
câu hỏi trung tâm là: đặc điểm tách từ của tiếng Việt ảnh hưởng thế nào đến việc chọn giữa hai họ mô hình?

Hai mô hình được fine-tune trong cùng điều kiện; epoch được chọn trên một tập dev tách từ train, và
validation (\NVal{} câu) được chấm đúng một lần. Để đưa PhoBERT --- vốn không có offset ký tự --- vào bài toán
trích xuất, nhóm cài một lớp tokenizer dựng offset từ bước tách từ \code{pyvi}. Ngoài EM/F1, nhóm đo: trần
EM do biên từ (không cần mô hình), phân loại lỗi tự động, precision/recall của hành vi từ chối, và kiểm định
cặp McNemar/bootstrap.

\textbf{Kết quả.} Điểm tổng của hai mô hình không khác nhau có ý nghĩa (EM \PhobertEM{} so với \XlmrEM; chênh
lệch \PairDiff, CI 95\% [\PairLo; \PairHi]), nhưng che hai hồ sơ ngược nhau: PhoBERT hơn \PairAnsDiff{} điểm EM
trên câu có đáp án, còn XLM-R hơn \,$-$\PairImpDiff{}\, điểm trên câu không có đáp án. Giả thuyết ban đầu
(PhoBERT yếu vì lỗi tách từ) \emph{sai} ở dạng tổng quát: chỉ \SegMisPct\% đáp án cắt ngang một từ, nên hiệu
ứng này bị chặn ở \SegCeilingLoss{} điểm EM, và phần lớn các trường hợp đó là lỗi của bộ tách từ. Điều quyết
định là hành vi từ chối: PhoBERT chỉ từ chối \PhobertAbsRate\% số câu, vì loss trên lớp ``không có đáp án''
của nó tăng khoảng 7 lần sau epoch 1 trong khi loss trên vị trí đáp án vẫn giảm --- hiện tượng lặp lại ở hai
tốc độ học. Cuối cùng, kiểm định bộ stress-test 250 câu của nhóm cho thấy chỉ \AuditAnsInCtx/\AuditAns{} câu có
đáp án chứa đáp án trong ngữ cảnh, và hệ thống luôn từ chối đứng đầu bảng điểm trên bộ này; bộ được giữ làm
đối tượng phân tích, không dùng để xếp hạng.

Mọi con số trong báo cáo được sinh tự động từ tệp \code{results/*.json} kèm commit.

\vspace{0.6em}
\noindent\textbf{Từ khoá:} đọc hiểu máy, hỏi đáp trích xuất, tiếng Việt, tách từ, PhoBERT, XLM-R,
UIT-ViQuAD~2.0, câu hỏi không có đáp án, đánh giá chẩn đoán.
